% ═══════════════════════════════════════════════════════════════
% ML-02b: Musterlösung Onboarding Widerstände (Reihenschaltung)
% Physik Klasse 8 | Doppelstunde 02b
% ═══════════════════════════════════════════════════════════════

\documentclass[11pt, a4paper]{article}

\usepackage[utf8]{inputenc}
\usepackage[T1]{fontenc}
\usepackage[ngerman]{babel}

\usepackage{helvet}
\renewcommand{\familydefault}{\sfdefault}

\usepackage{microtype}
\usepackage[margin=1.5cm, top=1.8cm, bottom=1.5cm]{geometry}
\usepackage{enumitem}
\usepackage{tabularx}
\usepackage{booktabs}
\usepackage{array}
\usepackage{tikz}
\usetikzlibrary{calc}

\usepackage{amsmath, amssymb}
\usepackage{siunitx}
\sisetup{locale = DE, per-mode = symbol}

\usepackage{fancyhdr}
\usepackage{lastpage}
\usepackage{needspace}

\usepackage{xcolor}
\usepackage{tcolorbox}
\tcbuselibrary{skins, breakable}

% Farben
\definecolor{physik}{HTML}{2c5aa0}
\definecolor{hellgrau}{HTML}{F5F5F5}
\definecolor{mittelgrau}{HTML}{888888}
\definecolor{dunkelgrau}{HTML}{444444}
\definecolor{rahmen}{HTML}{CCCCCC}
\definecolor{positive}{HTML}{c62828}
\definecolor{negative}{HTML}{1565c0}
\definecolor{loesung}{HTML}{c62828}

% Variablen
\newcommand{\fach}{Physik}
\newcommand{\thema}{Widerstände – Reihen- und Parallelschaltung}
\newcommand{\klasse}{8}

\colorlet{fachfarbe}{physik}

\newif\ifzeigesternchen
\zeigesternchentrue

% Kopf- und Fußzeile
\pagestyle{fancy}
\fancyhf{}
\renewcommand{\headrulewidth}{0.4pt}
\lhead{\small \fach{} Klasse \klasse{} | \thema{} – \textcolor{loesung}{\textbf{MUSTERLÖSUNG}}}
\rhead{\small Name: \underline{\hspace{3.5cm}}}
\cfoot{\small\textcolor{mittelgrau}{Seite \thepage{} von \pageref{LastPage}}}

% Boxen
\newtcolorbox{merkkasten}[1][]{
    colback=hellgrau,
    colframe=hellgrau,
    boxrule=0pt,
    arc=0pt,
    left=8pt, right=8pt, top=8pt, bottom=6pt,
    borderline north={2pt}{0pt}{fachfarbe},
    before skip=8pt, after skip=8pt,
    #1
}

\newtcolorbox{formelkasten}{
    colback=white,
    colframe=fachfarbe,
    boxrule=1.5pt,
    arc=0pt,
    left=8pt, right=8pt, top=6pt, bottom=6pt,
    before skip=6pt, after skip=6pt
}

\newtcolorbox{hinweiskasten}{
    colback=white,
    colframe=white,
    boxrule=0pt,
    arc=0pt,
    left=8pt, right=4pt, top=4pt, bottom=4pt,
    borderline west={3pt}{0pt}{fachfarbe}
}

% Befehle
\newcommand{\luecke}[1]{\underline{\hspace{#1}}}
\newcommand{\lsg}[1]{\textcolor{loesung}{\textbf{#1}}}
\newcommand{\punktebox}[1]{\hfill\small\textcolor{mittelgrau}{[\hspace{0.8em}/#1]}}
\newcommand{\stern}[1]{\ifzeigesternchen\textsuperscript{\textcolor{fachfarbe}{(#1)}}\fi}

\newcommand{\aufgabe}[2]{%
    \needspace{4\baselineskip}%
    \vspace{8pt}%
    \noindent\textbf{\textcolor{fachfarbe}{Aufgabe #1}} #2%
    \vspace{4pt}%
    \nopagebreak[4]%
}

\newcommand{\operator}[1]{\textbf{#1}}

\newlist{teil}{enumerate}{2}
\setlist[teil,1]{label=\alph*), leftmargin=1.8em, itemsep=4pt, parsep=0pt, topsep=4pt}

\newcommand{\antwortzeilen}[1]{%
    \par\vspace{4pt}%
    \foreach \n in {1,...,#1}{%
        \noindent\rule{\linewidth}{0.3pt}\vspace{8pt}%
    }%
}

\setlength{\parindent}{0pt}
\setlength{\parskip}{3pt}

\begin{document}

% ═══════════════════════════════════════════════════════════════
% HEADER
% ═══════════════════════════════════════════════════════════════

\begin{center}
    {\Large\bfseries\textcolor{fachfarbe}{\thema}}\\[0.2em]
    \textcolor{loesung}{\large\textbf{--- MUSTERLÖSUNG ---}}
\end{center}
\vspace{-0.5em}
\hrule
\vspace{0.8em}

% ═══════════════════════════════════════════════════════════════
% KASTEN 1: WIEDERHOLUNG OHMSCHES GESETZ
% ═══════════════════════════════════════════════════════════════

\begin{merkkasten}
\textbf{Kasten 1: Wiederholung – Das Ohmsche Gesetz}

\vspace{0.3em}
\begin{minipage}[t]{0.55\linewidth}
\textbf{Formel:} \quad $R = \dfrac{U}{I}$

\vspace{0.5em}
\textbf{Die drei Formen:}
\begin{itemize}[leftmargin=1.5em, itemsep=2pt]
    \item $R = $ \lsg{U / I} (Widerstand)
    \item $U = $ \lsg{R $\cdot$ I} (Spannung)
    \item $I = $ \lsg{U / R} (Stromstärke)
\end{itemize}

\vspace{0.3em}
\textbf{Einheit:} $[R] = 1\,\Omega = 1\,\dfrac{\text{V}}{\text{A}}$
\end{minipage}%
\begin{minipage}[t]{0.45\linewidth}
\begin{center}
\textbf{Formeldreieck:}

\vspace{0.3em}
\begin{tikzpicture}[scale=0.5]
    \draw[thick, fachfarbe] (0,0) -- (4,0) -- (2,3) -- cycle;
    \draw[fachfarbe] (0.5,1.5) -- (3.5,1.5);
    \node at (2,2.2) {\large $U$};
    \node at (0.9,0.6) {\large $R$};
    \node at (3.1,0.6) {\large $I$};
    \node at (2,0.6) {\large $\cdot$};
\end{tikzpicture}
\end{center}
\end{minipage}
\end{merkkasten}

% ═══════════════════════════════════════════════════════════════
% AUFGABE 1: KOPFRECHNEN
% ═══════════════════════════════════════════════════════════════

\aufgabe{1}{\stern{*} – Kopfrechnen mit dem Ohmschen Gesetz} \punktebox{8}

\operator{Berechne} im Kopf. \operator{Trage} die fehlenden Werte ein.

\vspace{0.5em}
\begin{center}
\begin{tabular}{|c|c|c|}
\hline
\textbf{Spannung $U$} & \textbf{Stromstärke $I$} & \textbf{Widerstand $R$} \\
\hline
12 V & 3 A & \lsg{4 $\Omega$} \\[0.5em]
\hline
\lsg{10 V} & 2 A & 5 $\Omega$ \\[0.5em]
\hline
24 V & \lsg{3 A} & 8 $\Omega$ \\[0.5em]
\hline
6 V & 0,5 A & \lsg{12 $\Omega$} \\[0.5em]
\hline
\end{tabular}
\end{center}

% ═══════════════════════════════════════════════════════════════
% KASTEN 2: REIHENSCHALTUNG
% ═══════════════════════════════════════════════════════════════

\begin{merkkasten}
\textbf{Kasten 2: Widerstände in Reihenschaltung}

\vspace{0.3em}
\begin{minipage}[t]{0.55\linewidth}
Wenn zwei Widerstände \textbf{hintereinander} geschaltet sind, addieren sich ihre Werte:

\begin{formelkasten}
\begin{center}
\large $R_{\text{ges}} = R_1 + R_2$
\end{center}
\end{formelkasten}

\vspace{0.3em}
\textbf{Wassermodell:} Zwei Engstellen hintereinander $\rightarrow$ doppelt so schwer durchzukommen!
\end{minipage}%
\begin{minipage}[t]{0.45\linewidth}
\begin{center}
\textbf{Schaltplan:}

\vspace{0.3em}
\begin{tikzpicture}[scale=0.45]
    % Geschlossener Stromkreis
    % Obere Leitung
    \draw[thick] (0,2) -- (1.5,2);
    \draw[thick] (3,2) -- (4,2);
    \draw[thick] (5.5,2) -- (7,2);
    % Rechte Leitung
    \draw[thick] (7,2) -- (7,0);
    % Untere Leitung
    \draw[thick] (7,0) -- (0,0);
    % Linke Leitung (mit Spannungsquelle)
    \draw[thick] (0,0) -- (0,0.7);
    \draw[thick] (0,1.3) -- (0,2);

    % Spannungsquelle (DIN: lange Linie = +, kurze dicke = -)
    \draw[positive, thick] (-0.4,1.3) -- (0.4,1.3); % + oben
    \draw[negative, line width=1.5pt] (-0.25,0.7) -- (0.25,0.7); % - unten
    \node[left] at (-0.5,1) {\tiny $U$};

    % R1 (Rechteck)
    \draw[thick, fachfarbe, fill=white] (1.5,1.7) rectangle (3,2.3);
    \node at (2.25,2) {\scriptsize $R_1$};

    % R2 (Rechteck)
    \draw[thick, fachfarbe, fill=white] (4,1.7) rectangle (5.5,2.3);
    \node at (4.75,2) {\scriptsize $R_2$};

    % Stromrichtung
    \draw[->, thick, positive] (3.5,2.5) -- (4.5,2.5);
    \node[above] at (4,2.5) {\tiny $I$};
\end{tikzpicture}
\end{center}
\end{minipage}
\end{merkkasten}

% ═══════════════════════════════════════════════════════════════
% AUFGABE 2: REIHENSCHALTUNG BERECHNEN
% ═══════════════════════════════════════════════════════════════

\aufgabe{2}{\stern{*}/\stern{**} – Gesamtwiderstand berechnen} \punktebox{12}

\begin{teil}
    \item \stern{*} $R_1 = 10\,\Omega$, $R_2 = 20\,\Omega$. \operator{Berechne} $R_{\text{ges}}$. \hfill (2 BE)

    \vspace{0.3em}
    $R_{\text{ges}} = $ \lsg{$10\,\Omega + 20\,\Omega = 30\,\Omega$}

    \vspace{0.5em}
    \item \stern{*} $R_1 = 15\,\Omega$, $R_2 = 25\,\Omega$. \operator{Berechne} $R_{\text{ges}}$. \hfill (2 BE)

    \vspace{0.3em}
    $R_{\text{ges}} = $ \lsg{$15\,\Omega + 25\,\Omega = 40\,\Omega$}

    \vspace{0.5em}
    \item \stern{**} $R_{\text{ges}} = 50\,\Omega$, $R_1 = 20\,\Omega$. \operator{Berechne} $R_2$. \hfill (3 BE)

    \vspace{0.3em}
    Umstellen: $R_2 = $ \lsg{$R_{\text{ges}} - R_1$} \hspace{0.5cm} Rechnung: \lsg{$50\,\Omega - 20\,\Omega = 30\,\Omega$}

    \vspace{0.5em}
    \item \stern{**} Eine Reihenschaltung ($R_1 = 30\,\Omega$, $R_2 = 20\,\Omega$) liegt an $U = 10\,\text{V}$. \hfill (5 BE)

    \operator{Berechne} $R_{\text{ges}}$ und $I$.

    \vspace{0.3em}
    $R_{\text{ges}} = $ \lsg{$30\,\Omega + 20\,\Omega = 50\,\Omega$}

    \vspace{0.3em}
    $I = \dfrac{U}{R_{\text{ges}}} = $ \lsg{$\dfrac{10\,\text{V}}{50\,\Omega} = 0{,}2\,\text{A}$}
\end{teil}

% ═══════════════════════════════════════════════════════════════
% AUFGABE 3: KENNLINIEN VERGLEICHEN
% ═══════════════════════════════════════════════════════════════

\aufgabe{3}{\stern{**} – Kennlinien vergleichen} \punktebox{5}

Die Simulation zeigt zwei Widerstände: einen ohmschen Widerstand (gerade Kennlinie) und eine Glühlampe (gekrümmte Kennlinie).

\begin{teil}
    \item \operator{Erkläre} mit einem Satz, warum der ohmsche Widerstand eine Gerade ergibt. \hfill (2 BE)

    \lsg{Bei einem ohmschen Widerstand ist R konstant, also ist I direkt proportional zu U ($I \sim U$), was eine Gerade durch den Ursprung ergibt.}

    \vspace{0.3em}
    \item \operator{Erkläre}, warum der Widerstand der Glühlampe bei höherer Spannung steigt. \hfill (3 BE)

    \textit{Hinweis: Denke an die Atomrümpfe im heißen Draht!}

    \lsg{Bei höherer Spannung fließt mehr Strom, der Draht wird heißer. Die Atomrümpfe schwingen stärker und behindern die Elektronen mehr. Dadurch steigt der Widerstand.}
\end{teil}

% ═══════════════════════════════════════════════════════════════
% KASTEN 3: PARALLELSCHALTUNG
% ═══════════════════════════════════════════════════════════════

\begin{merkkasten}
\textbf{Kasten 3: Widerstände in Parallelschaltung}

\vspace{0.3em}
\begin{minipage}[t]{0.55\linewidth}
Wenn zwei Widerstände \textbf{nebeneinander} geschaltet sind, wird der Gesamtwiderstand \textbf{kleiner}:

\begin{formelkasten}
\begin{center}
\large $\dfrac{1}{R_{\text{ges}}} = \dfrac{1}{R_1} + \dfrac{1}{R_2}$
\end{center}
\end{formelkasten}

\vspace{0.2em}
\textbf{Kurzformel:} $R_{\text{ges}} = \dfrac{R_1 \cdot R_2}{R_1 + R_2}$

\vspace{0.3em}
\textbf{Wassermodell:} Zwei Rohre nebeneinander $\rightarrow$ Wasser kann durch beide fließen!
\end{minipage}%
\begin{minipage}[t]{0.45\linewidth}
\begin{center}
\textbf{Schaltplan:}

\vspace{0.3em}
\begin{tikzpicture}[scale=0.45]
    % Geschlossener Stromkreis mit Parallelschaltung
    % Linke Leitung (mit Spannungsquelle)
    \draw[thick] (0,0) -- (0,0.7);
    \draw[thick] (0,1.3) -- (0,3);

    % Spannungsquelle (DIN)
    \draw[positive, thick] (-0.4,1.3) -- (0.4,1.3);
    \draw[negative, line width=1.5pt] (-0.25,0.7) -- (0.25,0.7);
    \node[left] at (-0.5,1) {\tiny $U$};

    % Obere Leitung zum Knotenpunkt
    \draw[thick] (0,3) -- (2,3);
    % Untere Leitung zum Knotenpunkt
    \draw[thick] (0,0) -- (2,0);

    % Verzweigung oben
    \draw[thick] (2,3) -- (2,2.3);
    \draw[thick] (2,3) -- (2,3.7);
    \draw[thick] (2,3.7) -- (2.5,3.7);
    \draw[thick] (2,2.3) -- (2.5,2.3);

    % R1 oben
    \draw[thick, fachfarbe, fill=white] (2.5,3.4) rectangle (4.5,4);
    \node at (3.5,3.7) {\scriptsize $R_1$};

    % R2 unten
    \draw[thick, fachfarbe, fill=white] (2.5,2) rectangle (4.5,2.6);
    \node at (3.5,2.3) {\scriptsize $R_2$};

    % Verbindung nach rechts
    \draw[thick] (4.5,3.7) -- (5,3.7);
    \draw[thick] (4.5,2.3) -- (5,2.3);
    \draw[thick] (5,3.7) -- (5,3);
    \draw[thick] (5,2.3) -- (5,3);

    % Rechte Leitung
    \draw[thick] (5,3) -- (6,3);
    \draw[thick] (6,3) -- (6,0);
    \draw[thick] (6,0) -- (2,0);

    % Stromrichtung
    \draw[->, thick, positive] (0.5,3.3) -- (1.5,3.3);
    \node[above] at (1,3.3) {\tiny $I$};
\end{tikzpicture}
\end{center}
\end{minipage}
\end{merkkasten}

% ═══════════════════════════════════════════════════════════════
% AUFGABE 4: PARALLELSCHALTUNG BERECHNEN
% ═══════════════════════════════════════════════════════════════

\aufgabe{4}{\stern{**}/\stern{***} – Parallelschaltung berechnen} \punktebox{15}

\begin{teil}
    \item \stern{**} $R_1 = 20\,\Omega$, $R_2 = 20\,\Omega$ parallel. \operator{Berechne} $R_{\text{ges}}$. \hfill (3 BE)

    \textit{Hinweis: Zwei gleiche Widerstände parallel $\rightarrow$ R\textsubscript{ges} = R/2}

    \vspace{0.3em}
    $R_{\text{ges}} = $ \lsg{$\dfrac{20\,\Omega}{2} = 10\,\Omega$}

    \vspace{0.5em}
    \item \stern{**} $R_1 = 30\,\Omega$, $R_2 = 60\,\Omega$ parallel. \operator{Berechne} $R_{\text{ges}}$. \hfill (4 BE)

    \vspace{0.3em}
    $R_{\text{ges}} = \dfrac{R_1 \cdot R_2}{R_1 + R_2} = $ \lsg{$\dfrac{30 \cdot 60}{30 + 60} = \dfrac{1800}{90} = 20\,\Omega$}

    \vspace{0.5em}
    \item \stern{***} \operator{Analysiere} die folgende Schaltung und \operator{berechne} $R_{\text{ges}}$. \hfill (8 BE)

    \vspace{0.5em}
    \begin{center}
    \begin{tikzpicture}[scale=0.5]
        % Gemischte Schaltung: R1 in Reihe mit (R2 || R3)
        % Linke Leitung
        \draw[thick] (0,0) -- (0,2);
        \draw[thick] (0,2) -- (1,2);

        % R1 in Reihe
        \draw[thick, fachfarbe, fill=white] (1,1.7) rectangle (3,2.3);
        \node at (2,2) {\small $R_1$};
        \node[below] at (2,1.6) {\scriptsize $10\,\Omega$};

        % Verbindung zur Parallelschaltung
        \draw[thick] (3,2) -- (4,2);

        % Verzweigung
        \draw[thick] (4,2) -- (4,3);
        \draw[thick] (4,2) -- (4,1);
        \draw[thick] (4,3) -- (4.5,3);
        \draw[thick] (4,1) -- (4.5,1);

        % R2 oben
        \draw[thick, fachfarbe, fill=white] (4.5,2.7) rectangle (6.5,3.3);
        \node at (5.5,3) {\small $R_2$};
        \node[above] at (5.5,3.4) {\scriptsize $20\,\Omega$};

        % R3 unten
        \draw[thick, fachfarbe, fill=white] (4.5,0.7) rectangle (6.5,1.3);
        \node at (5.5,1) {\small $R_3$};
        \node[below] at (5.5,0.6) {\scriptsize $20\,\Omega$};

        % Zusammenführung
        \draw[thick] (6.5,3) -- (7,3);
        \draw[thick] (6.5,1) -- (7,1);
        \draw[thick] (7,3) -- (7,2);
        \draw[thick] (7,1) -- (7,2);

        % Rechte Leitung
        \draw[thick] (7,2) -- (8,2);
        \draw[thick] (8,2) -- (8,0);
        \draw[thick] (8,0) -- (0,0);

        % Spannungsquelle
        \draw[positive, thick] (-0.4,1.3) -- (0.4,1.3);
        \draw[negative, line width=1.5pt] (-0.25,0.7) -- (0.25,0.7);
        \node[left] at (-0.6,1) {\small $U$};
    \end{tikzpicture}
    \end{center}

    \vspace{0.3em}
    \textit{Tipp: Berechne zuerst $R_{23}$ (parallel), dann $R_{\text{ges}} = R_1 + R_{23}$.}

    \vspace{0.3em}
    \lsg{Schritt 1: $R_{23} = \dfrac{20 \cdot 20}{20 + 20} = \dfrac{400}{40} = 10\,\Omega$}

    \vspace{0.3em}
    \lsg{Schritt 2: $R_{\text{ges}} = R_1 + R_{23} = 10\,\Omega + 10\,\Omega = 20\,\Omega$}
\end{teil}

\end{document}
